% Section 4.3: Competing Interpretations
% Structure: Intro → DM Interpretation → MSP Hypothesis → Counter-Arguments

\section{Competing Interpretations}
\label{sec:4.3}

The spectral and morphological properties established in the preceding section admit two fundamentally different explanations.
The first attributes the GCE to dark matter annihilation in the Galactic halo, drawing on the signal's consistency with theoretical WIMP predictions.
The second invokes a large population of unresolved millisecond pulsars concentrated in the Galactic bulge, supported by photon-count statistics and morphological studies.
We review each interpretation in turn, followed by independent astrophysical constraints that challenge the pulsar hypothesis.

\subsection{The Dark Matter Interpretation}
\label{sec:4.3.1}

The properties of the Galactic Center Excess align closely with the predictions of WIMP dark matter annihilation, a coincidence that has sustained interest in this interpretation for over a decade.
The strength of the dark matter case rests on three independent lines of agreement---spectral, morphological, and normalization---each of which is tightly constrained by a small number of parameters.

The expected gamma-ray flux from dark matter annihilation in the Galactic halo takes the factorized form
\begin{equation}
	\frac{d\Phi_\gamma}{d\Omega\, dE} = \frac{\langle \sigma v \rangle}{8\pi\, m_\chi^2}\, \frac{dN_\gamma}{dE}\, \times \int_{\mathrm{l.o.s.}} \rho^2(r)\, dl\,,
	\label{eq:dm_flux_gce}
\end{equation}
where the particle physics content---the annihilation cross section $\langle \sigma v \rangle$, the WIMP mass $m_\chi$, and the photon yield per annihilation $dN_\gamma/dE$---is cleanly separated from the astrophysical J-factor, the line-of-sight integral of the squared dark matter density (see Section~\ref{sec:1.4} for a detailed treatment).
Because annihilation occurs essentially at rest, the observed spectrum is fixed by the mass and annihilation channel, while the angular distribution on the sky is determined entirely by the density profile.
With only three free parameters---$m_\chi$, the annihilation channel, and the inner slope $\gamma$ of the halo profile---the dark matter model must simultaneously reproduce both the energy spectrum and the spatial morphology of the excess~\cite{Daylan:2014rsa}.

On the spectral side, the GCE's peaked emission at $1$--$3$~GeV with a steep sub-GeV rise and a power-law decline above the peak is precisely what one expects from the hadronization of $b\bar{b}$ final states produced by a WIMP in the $\sim 30$--$50$~GeV mass range.
Multiple independent analyses have converged on consistent best-fit masses: Daylan et al.
\ found $36$--$51$~GeV~\cite{Daylan:2014rsa}, Abazajian et al.\ obtained $39.4^{+3.7}_{-2.9}$~GeV (statistical)~\cite{Abazajian:2014fta}, and Calore, Cholis \& Weniger reported $49^{+6.4}_{-5.4}$~GeV after marginalizing over 60 Galactic diffuse emission models~\cite{Calore:2014xka}.
This spectral agreement extends beyond the $b\bar{b}$ channel: annihilation into $c\bar{c}$ provides comparably good fits, while leptonic channels ($\tau^+\tau^-$) are disfavored at the 95\% confidence level~\cite{Calore:2014xka}.

The morphological agreement is equally compelling.
The excess is approximately spherically symmetric and centered within $\sim 0.05^\circ$ of Sgr~A$^*$, with axis ratios departing from unity by more than $\sim 20$\% yielding substantially worse fits~\cite{Daylan:2014rsa}.
The signal extends to at least $\sim 10^\circ$ ($\sim 1.5$~kpc) from the Galactic Center, with a radial intensity profile that falls off as $F_\gamma \propto r^{-(2.2\text{--}2.6)}$, implying a dark matter density scaling as $\rho \propto r^{-\gamma}$ with $\gamma \approx 1.1$--$1.3$~\cite{Daylan:2014rsa, Abazajian:2014fta, Calore:2014xka}.
This steepened inner slope, slightly above the canonical NFW value of $\gamma = 1.0$, is physically motivated: the gravitational influence of baryonic matter in the inner Galaxy can adiabatically contract the dark matter cusp, producing exactly such a steepening~\cite{Goodenough:2009gk}.
No known astrophysical emission mechanism naturally produces a spherically symmetric signal with this radial profile.
The observed emission does not trace any combination of gas, dust, star formation, or stellar populations---all of which are concentrated along the Galactic disk---but instead follows the square of the anticipated dark matter density~\cite{Daylan:2014rsa}.

A notable feature of the dark matter fit is its normalization.
The annihilation cross section required to produce the observed flux falls in the range $\langle \sigma v \rangle \sim (1$--$3) \times 10^{-26}$~cm$^3$/s~\cite{Daylan:2014rsa, Calore:2014xka}, which is consistent with the canonical thermal relic value of $\langle \sigma v \rangle_{\mathrm{cosmo}} \approx 2.2 \times 10^{-26}$~cm$^3$/s~\cite{Cirelli:2024ssz}.
In other words, the very cross section needed to explain the GCE brightness today is the same cross section that would cause a WIMP to freeze out in the early universe with an abundance equal to the observed cosmological dark matter density.
No Sommerfeld enhancement, substructure boost factors, or non-thermal production mechanisms are required~\cite{Daylan:2014rsa}.
The fact that three independent observables---spectrum, morphology, and flux normalization---converge on a simple, internally consistent dark matter model with parameters close to the thermal relic prediction has made the GCE one of the most compelling indirect detection signals in the literature.
It should be noted, however, that Fermi-LAT observations of dwarf spheroidal galaxies place upper limits on the annihilation cross section that are in tension with this interpretation for some channel and mass combinations~\cite{Fermi-LAT:2015att}, and the inferred parameters depend on assumptions about the local dark matter density and the shape of the halo profile at small radii.

\subsection{The Millisecond Pulsar Hypothesis}
\label{sec:4.3.2}

The leading astrophysical alternative to dark matter attributes the GCE to a large population of unresolved millisecond pulsars (MSPs) in the inner Galaxy.
First proposed by Abazajian in 2011~\cite{Abazajian:2010zy}, this hypothesis builds on the observation that MSPs produce gamma-ray spectra naturally similar to the GCE: a stacked analysis of 61 Fermi-LAT-detected MSPs yields a spectrum well described by a power law with spectral index $\alpha \approx 1.57$ and an exponential cutoff at $E_{\mathrm{cut}} \approx 3.78$~GeV, parameters that fall within the 99\% confidence contours of the GCE spectral fit~\cite{Calore:2014xka}.
Unlike the dark matter interpretation, which requires a new particle, the MSP hypothesis invokes a known class of astrophysical source---the question is whether a sufficiently large and concentrated population exists in the Galactic Center region.

The spatial distribution of these putative MSPs presents the first challenge.
The known gamma-ray MSP population follows a disk-like distribution, which would produce a signal elongated along the Galactic plane with an axis ratio of $\sim 1$:$6$---incompatible with the observed spherical morphology of the GCE~\cite{Daylan:2014rsa}.
A standard disk MSP population can account for at most 5--10\% of the excess flux~\cite{Calore:2014xka}.
To reproduce the concentrated, roughly spherical GCE morphology, the MSP number density would need to increase steeply toward the Galactic Center as approximately $r^{-2.5}$, requiring a distinct ``bulge'' population rather than the familiar disk population~\cite{Calore:2014xka}.
Such a bulge component is not unreasonable---the high stellar densities in the inner Galaxy could produce MSPs through the same dynamical channels operating in globular clusters---but it cannot easily account for the signal's extension to $\sim 10$--$15^\circ$ from the Galactic Center~\cite{Calore:2014xka}.

Two independent statistical analyses published in 2016 provided the strongest evidence favoring the point-source interpretation.
Lee et al.
\ developed the Non-Poissonian Template Fitting (NPTF) method, which generalizes standard template fitting to account for the photon-count statistics expected from unresolved point sources~\cite{Lee:2015fea}.
While a smooth dark matter halo produces Poissonian photon counts in each pixel (the variance equals the mean), a population of faint point sources introduces additional variance through localized photon clusters.
By modeling the source-count distribution $dN/dF$ as a broken power law, Lee et al.
\ found that point sources distributed according to an NFW-squared spatial template fully absorbed the GCE, with a Bayes factor of $\sim 10^6$ favoring the point-source model over smooth dark matter emission.
Their inferred source-count function rose sharply just below Fermi's detection threshold, with the break corresponding to a luminosity of $L_\gamma \sim 2$--$3 \times 10^{34}$~erg/s.
Extrapolating this distribution, they estimated that $\sim 400$ near-threshold sources could account for the entire excess within $10^\circ$ of the Galactic Center~\cite{Lee:2015fea}.

In a complementary approach, Bartels, Krishnamurthy \& Weniger applied a wavelet decomposition to the Fermi-LAT maps, which is sensitive to small-scale spatial fluctuations produced by unresolved point sources.
Their analysis independently detected excess power at small angular scales in the inner Galaxy, consistent with a population of sub-threshold sources contributing to the GCE~\cite{Bartels:2015aea}.

A separate line of evidence came from morphological studies.
Macias et al.
\ argued that the GCE spatial distribution is better described by the stellar mass of the Milky Way's boxy/X-shaped bulge than by a spherically symmetric NFW profile~\cite{Macias:2016nev, Macias:2019omb}.
Using near-infrared maps to trace the three-dimensional structure of the central stellar bar, they found that a bulge template absorbed the excess emission with a statistically better fit than the NFW template.
Because MSPs are expected to trace the old stellar population that forms the bulge, this morphological preference was interpreted as further evidence for an astrophysical origin.
Similar conclusions were reached by Bartels et al.
\ \cite{Bartels:2018xom}, who found that the GCE traces the stellar mass distribution in the inner Galaxy, suggesting an intimate connection between the excess and the Galactic bulge population.

Taken together, the NPTF photon statistics, wavelet power, and stellar bulge morphology built a case that appeared to decisively favor the MSP interpretation by the late 2010s.
As we will see in Section~\ref{sec:4.4}, however, subsequent work revealed that each of these lines of evidence is more fragile than initially appreciated.

\subsection{Counter-Arguments Against MSPs}
\label{sec:4.3.3}

Despite the statistical evidence from NPTF and wavelet analyses, several independent lines of reasoning have challenged the viability of the MSP interpretation.
These arguments do not rely on the photon statistics of the GCE itself, but instead constrain the expected properties of any MSP population large enough to produce the excess.

The first constraint comes from scaling the observed relationship between low-mass X-ray binaries (LMXBs) and MSPs.
Because MSPs evolve directly from LMXBs, stellar populations of similar age should contain a predictable ratio between the two classes.
Cholis, Hooper \& Linden exploited this connection by measuring the gamma-ray output of 16 Fermi-detected globular clusters hosting a total of 5 bright ($L > 10^{36}$~erg/s) LMXBs, finding that their combined luminosity amounts to only 4.8\% of the GCE flux within the inner $5^\circ$~\cite{Cholis:2014lta}.
If MSPs were responsible for the entire excess, the inner Galaxy should contain $\sim 100$ bright LMXBs; INTEGRAL observations have detected only 7 candidates in this region.
Even without correcting for selection biases that inflate the globular cluster gamma-ray yield, this comparison limits the MSP contribution to $\sim 4$--$11$\% of the GCE.
After accounting for these biases, the authors estimated that only $\sim 1$--$5$\% of the excess can plausibly originate from unresolved MSPs~\cite{Cholis:2014lta}.

A second challenge arises from the dynamics of MSP formation.
Neutron stars receive natal velocity kicks of order $\sim 10^2$~km/s during their formation, and while the kicks imparted to recycled pulsars are typically smaller than those of young isolated pulsars, they remain large enough to significantly broaden the spatial distribution of the resulting MSP population over gigayear timescales~\cite{Brandt:2015ula}.
Brandt \& Kocsis showed that this kinetic injection should disperse an in-situ MSP population well beyond the concentrated profile inferred for the GCE, particularly in the innermost $\sim 1^\circ$ ($\sim 150$~pc) where the observed emission requires a steep volumetric density scaling as $\rho \propto r^{-2.4}$ to $r^{-2.6}$~\cite{Brandt:2015ula}.
A kicked population would produce a significantly flatter, more extended profile, inconsistent with the observed concentration of the signal.

The third and perhaps most direct argument relies on the luminosity function of known MSP populations.
If the pulsars responsible for the GCE share the luminosity function of MSPs observed in either the Galactic plane or in globular clusters, then the bright tail of this distribution should extend above Fermi's detection threshold, producing individually resolved sources.
Hooper \& Linden first quantified this constraint using the globular cluster MSP luminosity function, finding that Fermi should have detected a significant number of individual pulsars from the GCE population~\cite{Hooper:2016rap}.
The analysis was recently updated by Amerio, Hooper \& Linden using 15.8 years of Fermi data and 157 globular clusters, yielding a mean MSP luminosity of $\langle L_\gamma \rangle \sim (1$--$8) \times 10^{33}$~erg/s with a log-normal width $\sigma_L \sim 1.4$--$2.8$~\cite{Amerio:2024qor}.
Under this luminosity function, Fermi's source catalogs should contain $N_{\mathrm{MSP}} \sim 17$--$37$ individually resolved pulsars from the GCE population.
Similarly, adopting the luminosity function of Galactic plane MSPs predicts $\sim 10$--$35$ detections~\cite{Holst:2024fvb}.
In reality, only three millisecond pulsar candidates have been identified with directions and distances compatible with an inner Galaxy population~\cite{Amerio:2024qor, Holst:2024fvb}.
The discrepancy between the $\sim 10$--$37$ expected detections and the 3 observed candidates implies that the GCE can originate from MSPs only if those pulsars are systematically less luminous than any known MSP population---a requirement for which no astrophysical mechanism has been identified.
This ``missing bright pulsars'' problem represents one of the strongest empirical challenges to the MSP hypothesis, and it is the constraint that the analysis presented later in this chapter (Section~\ref{sec:4.6}) directly addresses.
